% --- PREÁMBULO FINAL OPTIMIZADO PARA LA EDITORIAL UD ---

% 1. CLASE DEL DOCUMENTO Y OPCIONES GLOBALES (PRIORIDAD DE LA PLANTILLA)
% NOTA: La clase "book" es la base, las opciones son las del formato UD.
%\documentclass[12pt,twoside,spanish]{book}

% 2. CONFIGURACIÓN DEL LENGUAJE Y TIPOGRAFÍA UD (PRIORIDAD AL DISEÑO)

\usepackage{fontspec} % Necesario para LuaLaTeX y la tipografía UD
\input{esConfigDir/esConfigFonts} % Definición de fuentes (lhseries, medseries, etc.)


% 3. TIPOGRAFÍA Y ESTILOS DE LA EDITORIAL UD (FUENTES Y GEOMETRÍA)
% NOTA: Se elimina tu \usepackage{geometry} y tu \usepackage{fontspec} manual
% para usar la configuración exacta de la editorial.


\usepackage{geometry}
\geometry{
	paperheight=24.0cm,
	paperwidth=17.0cm,
	textwidth=12.5cm,
	tmargin=2.5cm,
	bmargin=2.5cm,
	left=2.5cm,
	right=2.0cm
	%   bindingoffset=1.1cm % Usar si se desea ajuste al encuadernar
}

% 4. ESTILOS GRÁFICOS Y COMANDOS UD (\portadillauno, \espart, \nota, \fuente, 'esstyle')

%%% Editar %%%
\def\texttitulocorto{Título corto del libro}
%Lista de autores
\def\loauthors{Nombre o nombres de los autores del libro separados por comas}
%%%%%%%%%%%%%%

% Estos son necesarios para las portadillas y las cornisas especiales.
%% Estilo LaTeX Espacios Editorial UD
%% Autores: Editores UDFJC
%% Versión:  v0.1 30MAR2024
%% Compilar con LuaLaTeX, i.e: lualatex main.tex
\usepackage[spanish]{babel}
\usepackage{csquotes}
\usepackage{caption}
\usepackage{titlesec}
\usepackage{titletoc}
\usepackage[titles]{tocloft}

%\setlength{\parindent}{1.25cm}
%\setlength\parindent{0pt}

\usepackage{xcolor}
\definecolor{MSBlack}{rgb}{0.0, 0.0, 0.0}
%\definecolor{MSBLack}{rgb}{.16,.16,.16}
\definecolor{MSLightBlack}{rgb}{.10,.10,.10}
\definecolor{MSCustomGray}{RGB}{144, 147, 150}
\definecolor{MSCustomBlackT}{RGB}{20, 20, 20}
\definecolor{MSCustomBlackST}{RGB}{58, 59, 66}

\newcommand\fontstit{\medseries}
\newcommand\fontschap{\regseries}
\newcommand\fontssec{\regseries}
\newcommand\fontssubsec{\regseries}
\newcommand\fontssubsubsec{\regseries}
\newcommand\fontsparagraph{\regseries\itshape}

% Small caps
\newcommand{\smallcaps}[1]{\textls[25]{\textsc{\MakeLowercase{#1}}}}

%configuracion para titulos de partes, capitulos, secciones, subsecciones, subsubsecciones
\newcommand{\chapterlabel}{\regseries Capítulo~\thechapter.}

% \titleformat{<command>}[<shape>]{<format>}{<label>}{<sep>}{<before-code>}[<after-code>]

\titleformat{\part}[display]{\color{MSBlack}\fontsize{25pt}{2.0mm}\selectfont\fontstit}{\thepart}{12pt}{\raggedright}
\titleformat{\chapter}[block]{\color{MSBlack}\fontsize{25pt}{2.5mm}\selectfont\fontschap}{\chapterlabel}{0.3em}{\raggedright}
\titleformat{\section}[block]{\color{MSBlack}\fontsize{16pt}{2.0mm}\selectfont\fontssec}{\thesection}{0.5em}{\raggedright}
\titleformat{\subsection}[block]{\color{MSBlack}\fontsize{13pt}{2.0mm}\selectfont\fontssubsec}{\thesubsection}{0.5em}{\raggedright}
\titleformat{\subsubsection}[block]{\color{black}\fontsize{12pt}{2.0mm}\selectfont\fontssubsubsec}{\thesubsubsection}{0.5em}{\raggedright}
\titleformat{\paragraph}[block]{\color{black}\fontsize{12pt}{2.0mm}\selectfont\fontsparagraph}{\theparagraph}{0.5em}{\raggedright}

%Espacios para titulos
\titlespacing{\part}{0cm}{10cm}{2.8cm}
\titlespacing{\chapter}{0cm}{1cm}{2.8cm}
\titlespacing{\section}{0cm}{7mm}{2mm}
\titlespacing{\subsection}{0cm}{5mm}{2mm} 
\titlespacing{\subsubsection}{0cm}{5mm}{0mm}
\titlespacing{\paragraph}{0cm}{5mm}{0mm}

% Configuración para apéndices
\makeatletter
\g@addto@macro{\appendix}{%
\titleformat{\chapter}[block]{\color{MSBlack}\fontsize{25pt}{2.5mm}\selectfont\fontschap}{Apéndice~\thechapter.}{0.3em}{\raggedright}
  \titlecontents{chapter}[0em]
    {\vspace{1.5mm}\medseries\normalsize}
    {Apéndice~\thecontentslabel.\enspace}
    {}
    {\hfill\small\lhseries\contentspage\vspace{2.5mm}}[]
}
\makeatother

%configuracion tabla de contenido
\titlecontents{part}[0em]
{\bf \sbseries\normalsize}%
{}% numbered sections formatting
{}% unnumbered sections formatting
{\titlerule*[1pc]{}\small\lhseries\contentspage \vspace{1.0mm}}[]%

\titlecontents{chapter}[0em]
{\vspace{1.5mm}\medseries\normalsize}%
{Capítulo~\thecontentslabel.\enspace}% numbered sections formatting
{}% unnumbered sections formatting
{\titlerule*[1pc]{}\footnotesize\lhseries\contentspage \vspace{1mm}}[]%

\titlecontents{section}[1.5em]
{\small\lhseries}%
{}% numbered sections formatting
{}% unnumbered sections formatting
{\titlerule*[1pc]{}\footnotesize\lhseries\contentspage \vspace{0.5mm}}[]%

\titlecontents{subsection}[3em]
{\lhseries}%
{}% numbered sections formatting
{}% unnumbered sections formatting
{\titlerule*[1pc]{}\small\lhseries\contentspage \vspace{0.5mm}}[]%

\usepackage{setspace}
\linespread{1.1}
\usepackage[skip=3.5pt, indent=20pt]{parskip}

\usepackage{fancyvrb}
\usepackage{fancyhdr}

\fancypagestyle{esstyle}{%
\fancyhead{}
  \fancyfoot{}
  \renewcommand{\headrulewidth}{0pt}
%   \newcommand{\runningtitle}{\fontsparagraph \scriptsize \texttitulocorto}
%  \fancyhf[HLE]{\lhseries\scriptsize Author's Name}
  \fancyhf[HLE]{\lhseries\scriptsize \loauthors}
%  \fancyhf[HRO]{\runningtitle}
  \fancyhf[HRO]{\fontsparagraph \scriptsize \texttitulocorto}
  \fancyfoot[LE]{{\fontsize{10.5}{1}\usefont{T1}{Roboto-TLF}{thin}{}\selectfont{E}}\hspace{-0.5mm}{\footnotesize\usefont{T1}{Inter-Sup}{thin}{}\selectfont{\reflectbox{S}}}{\usefont{T1}{Roboto-TLF}{thin}{}\selectfont{\thinspace|\thinspace}}{\lhseries\scriptsize\thepage}}

  \fancyfoot[RO]{{\lhseries\scriptsize\thepage}{\usefont{T1}{Roboto-TLF}{thin}{}\selectfont{\thinspace|\thinspace}}{\fontsize{10.5}{1}\usefont{T1}{Roboto-TLF}{thin}{}\selectfont{E}}\hspace{-0.5mm}{\footnotesize\usefont{T1}{Inter-Sup}{thin}{}\selectfont{\reflectbox{S}}}}
  \fancyfoot[C]{}
}

%https://www.latex-project.org/help/documentation/fntguide.pdf

\fancypagestyle{plain}{%
  \fancyhead{}
  \fancyfoot{}
  \renewcommand{\headrulewidth}{0pt}
  \fancyfoot[LE]{{\fontsize{10.5}{1}\usefont{T1}{Roboto-TLF}{thin}{}\selectfont{E}}\hspace{-0.5mm}{\footnotesize\usefont{T1}{Inter-Sup}{thin}{}\selectfont{\reflectbox{S}}}{\usefont{T1}{Roboto-TLF}{thin}{}\selectfont{\thinspace|\thinspace}}{\lhseries\scriptsize\thepage}}

  \fancyfoot[RO]{{\lhseries\scriptsize\thepage}{\usefont{T1}{Roboto-TLF}{thin}{}\selectfont{\thinspace|\thinspace}}{\fontsize{10.5}{1}\usefont{T1}{Roboto-TLF}{thin}{}\selectfont{E}}\hspace{-0.5mm}{\footnotesize\usefont{T1}{Inter-Sup}{thin}{}\selectfont{\reflectbox{S}}}}
  \fancyfoot[C]{}
}

\assignpagestyle{\chapter}{plain}

%\setcounter{table}{0}
\counterwithout{table}{chapter}
\counterwithout{figure}{chapter}
\counterwithout{equation}{chapter}
\captionsetup[figure]{labelfont=bf, labelsep=period, font=footnotesize, justification=centering, width = 0.8\textwidth}
\captionsetup[table]{labelfont=bf, labelsep=period, font=footnotesize, justification=centering, width = 0.8\textwidth}

%footnote format
\makeatletter
\renewcommand{\@makefntext}[1]{%
  \begin{list}{}{%
    \setlength{\labelwidth}{1.5em}%
    \setlength{\leftmargin}{\labelwidth}%
    \setlength{\labelsep}{3pt}%
    \setlength{\itemsep}{0pt}%
    \setlength{\partopsep}{0pt}%
    \setlength{\topsep}{0pt}%
    \footnotesize}%
  \item[\@makefnmark\hfil]#1%
  \end{list}%
}
\makeatother

\newcommand{\nota}[1]{\parbox{10cm}{\centering\scriptsize\color{MSBlack} {\bf Nota: }{#1}}}
\newcommand{\fuente}[1]{\parbox{10cm}{\centering\scriptsize\color{MSBlack} {\bf Fuente: }{#1}}}

\usepackage{wrapfig}

\newcommand{\espart}[6]{\part[Parte #1]{\protect\thispagestyle{empty}
%{
%\hspace{-2.75cm}\centering\includegraphics[width=17cm, height=8cm]{parte1_img}
%\begin{figure}[!h] %this figure will be at the right
\vspace{-14.0cm}
   \begin{figure}[ht!]\hspace{-2.9cm}\includegraphics[width=1.045\paperwidth, height=0.5\paperheight]{#2}\end{figure}
%\end{figure}
\flushleft
  \medseries{\Large\noindent Parte {#3}.}
  ~\\
  ~\\
  ~\\
  ~\\
  \lhseries{#4}
  ~\\
  ~\\
  ~\\
  ~\\
  \lhseries{\Large #5}
  ~\\
  ~\\
  ~\\
  ~\\
  \lhseries{\normalsize #6}
}}

%%%%%%%%%%%%%%%%%%


%%%%%%%%%% PORTADILLA 1 %%%%%%%%%%%%%%%
\newcommand{\portadillauno}[2]{
\vspace*{6.5cm}
\begin{flushleft}
  \linespread{3}
%  \color{MSCustomBlackT}\lsstyle
  \color{MSBlack}\lsstyle
  \fontsize{22pt}{3.0mm} \lhseries

  \leftskip=2mm
  #1   
  \vspace{2mm}
  \leftskip=-5mm

  %{\color{white}\colorbox{MSCustomGray}
    %{\parbox[][1.5cm][c]{15.2cm}{
        %\parbox{12.5cm}{\leftskip=7mm \fontsize{17pt}{2.5mm} \lhseries \vspace{1.0mm}
%          Consed ut escimos num sequid ut aut mos {\break} excerro et et, tempos quam, optat fuga.\vspace{0.5mm}
          %#2 \vspace{0.5mm}
        %}
      %}
    %}
  %}

\end{flushleft}}
%%%%%%%%%% FIN PORTADILLA 1 %%%%%%%%%%%%%%%

%%%%%%%%%% PORTADILLA 2 %%%%%%%%%%%%%%%
\newcommand{\portadillados}[2]{
\vspace*{6.5cm}
\begin{flushleft}
  \linespread{3}
%  \color{MSCustomBlackT}\lsstyle
  \color{MSBlack}\lsstyle
  \fontsize{22pt}{3.0mm} \lhseries

  \leftskip=2mm
  #1
   
  \vspace{2mm}
  \leftskip=-5mm

  {\parbox[][1.5cm][c]{15.2cm}{
      \parbox{12.5cm}{\leftskip=7mm \fontsize{17pt}{2.5mm} \lhseries \vspace{1.0mm}
%          Consed ut escimos num sequid ut aut mos {\break} excerro et et, tempos quam, optat fuga.\vspace{0.5mm}
          #2 \vspace{0.5mm}
      }
    }
  }
\end{flushleft}}
%%%%%%%%%% PORTADILLA 2 %%%%%%%%%%%%%%%

\makeatletter
\newcommand*\reftoanex[1]{\def\@currentlabelname{#1}}
\makeatother

%% COLOFON ESTILOS
\definecolor{MSGreyCOLOFON}{rgb}{.35,.35,.35}
\definecolor{MSGreyLATEX}{rgb}{.55,.55,.55}

%\usepackage{libertine}
\DeclareRobustCommand\myLaTeX{%
    L\kern-.3em%
    \raisebox{.4ex}{\scshape a}\kern-.13em%
    T\kern-.2em%
    \raisebox{-.5ex}{E}\kern-.13em%
    X%
}

\newcommand{\escolofon}[2]{
\thispagestyle{empty}
\vspace*{10.0cm}
\begin{tikzpicture}[remember picture,overlay]
  \node[text=gray,rotate=0,anchor=east, yshift=5mm] at (15.5,11.5){
    \includegraphics[width=5.0cm, height=4.0cm]{esIcons/logoESColofon_corner.png}
  };
\end{tikzpicture}

\vspace{-4.5cm}
\begin{center}
{\fontsize{7pt}{7pt}\selectfont
\color{MSGreyCOLOFON}
Este libro fue editado y diagramado en} {\fontsize{13pt}{13pt}\selectfont\color{MSGreyLATEX}{\myLaTeX}}
\vspace{3.5cm}

\includegraphics[scale=0.05]{esIcons/logoESColofon.png}
  
  \vspace{0.5cm}
  {\scriptsize
    \color{MSGreyCOLOFON}Este libro se término de imprimir en #1 de #2\\ en la Editorial UD. Bogotá, Colombia.}
\end{center}
}
%%%%%%%%%%%%%%%%%%


%\usepackage{showframe}% Similar al calderon, apagar al terminar de maquetar.
\usepackage[axes,cam, width=215truemm, height=275truemm, lualatex,center]{crop} % Descomentar solo en producción final

% 5. AJUSTES FINALES DE FORMATO DEL PÁRRAFO
\usepackage{microtype} % Microtipografía avanzada
\usepackage{ragged2e} % Para control de alineación
\usepackage{authblk} % Manejo de autores
\decimalpoint % Para usar la coma como separador decimal.
\setlength\parindent{0pt} % Sin sangría al inicio de párrafos, como pide la UD.
\newcommand{\normalfonts}{\fontsize{12pt}{14.4pt}\selectfont}
\let\normalsize\normalfonts


%%%%%%%%%%%%%%%%%%%%%%%%%%%%%%%%%%%%%%%%%%%%%%%%%%%%%%%%%%%%%%%%%%%%%%%%%%%%%%%%%%%%%%%%%%%%%%%%%
%%% 6. BIBLIOGRAFÍA (NORMA IEEE - EDITORIAL UD)
%%%%%%%%%%%%%%%%%%%%%%%%%%%%%%%%%%%%%%%%%%%%%%%%%%%%%%%%%%%%%%%%%%%%%%%%%%%%%%%%%%%%%%%%%%%%%%%%%
\usepackage[hidelinks, unicode]{hyperref}
\urlstyle{same} 

\usepackage[
backend=biber,
style=ieee,      
sorting=none,    
isbn=false,      
doi=true         
]{biblatex}

\addbibresource{Referencias/references.bib}

% AJUSTE DE ESPACIO (Lo que solicitaste)
\setlength{\bibitemsep}{1.5ex} 

% CORRECCIONES PARA IDIOMA ESPAÑOL
\DefineBibliographyStrings{spanish}{
	and = {y},              
	bibliography = {Referencias},
	references = {Referencias},
	andothers = {et~al\adddot} 
}

\DeclareDelimFormat{finalnamedelim}{%
	\ifnumgreater{\value{liststop}}{2}{\finalandcomma\addspace}{}%
	\addspace\bibstring{y}\addspace}



% ÍNDICES Y SUBARCHIVOS
%\usepackage{makeidx}
%\usepackage{subfiles}

%%%%%%%%%%%%%%%%%%%%%%%%%%%%%%%%%%%%%%%%%%%%%%%%%%%%%%%%%%%%%%%%%%%%%%%%%%%%%%%%%%%%%%%%%%%%%%%%%
%%% 7. CORNISAS Y ESTILO DE PÁGINA (PRIORIDAD DE LA PLANTILLA Y TUS REQUERIMIENTOS)
%%%%%%%%%%%%%%%%%%%%%%%%%%%%%%%%%%%%%%%%%%%%%%%%%%%%%%%%%%%%%%%%%%%%%%%%%%%%%%%%%%%%%%%%%%%%%%%%%

\usepackage{emptypage}

% Definición del estilo 'plain' para páginas de inicio de capítulo (sin cornisas)
%\fancypagestyle{plain}{
	%\fancyhf{} % Limpiar todo
%\renewcommand{\headrulewidth}{0pt} % Quitar la línea
%\fancyfoot[C]{\thepage} % Número de página centrado en el pie (estándar)
%}
% Definición del estilo 'plain' - AHORA COINCIDE CON EL PIE DE esstyle
\fancypagestyle{plain}{%
	\fancyhf{} % Limpiar todo (incluyendo el número centrado anterior)
\renewcommand{\headrulewidth}{0pt} % Quitar la línea
	% Pie de página izquierdo (Páginas Pares, E.S. | #)
\fancyfoot[LE]{{\fontsize{10.5}{1}\usefont{T1}{Roboto-TLF}{thin}{}\selectfont{E}}\hspace{-0.5mm}{\footnotesize\usefont{T1}{Inter-Sup}{thin}{}\selectfont{\reflectbox{S}}}{\usefont{T1}{Roboto-TLF}{thin}{}\selectfont{\thinspace|\thinspace}}{\lhseries\scriptsize\thepage}}
% Pie de página derecho (Páginas Impares, # | E.S.)
\fancyfoot[RO]{{\lhseries\scriptsize\thepage}{\usefont{T1}{Roboto-TLF}{thin}{}\selectfont{\thinspace|\thinspace}}{\fontsize{10.5}{1}\usefont{T1}{Roboto-TLF}{thin}{}\selectfont{E}}\hspace{-0.5mm}{\footnotesize\usefont{T1}{Inter-Sup}{thin}{}\selectfont{\reflectbox{S}}}}
\fancyfoot[C]{} % Dejar esto vacío para asegurar que el centro está limpio
}



%Esto asegura que la página izquierda de un capítulo nuevo esté vacía (sin cornisas)
\renewcommand{\cleardoublepage}{%
\clearpage\ifodd\value{page}\else\hbox{}\thispagestyle{empty}\newpage\fi%
}

%%%%%%%%%%%%%%%%%%%%%%%%%%%%%%%%%%%%%%%%%%%%%%%%%%%%%%%%%%%%%%%%%%%%%%%%%%%%%%%%%%%%%%%%%%%%%%%%%
%%% 8. PAQUETES FUNCIONALES Y DE CONTENIDO (MANTENIDOS)
%%%%%%%%%%%%%%%%%%%%%%%%%%%%%%%%%%%%%%%%%%%%%%%%%%%%%%%%%%%%%%%%%%%%%%%%%%%%%%%%%%%%%%%%%%%%%%%%%

% PAQUETES DE MATEMÁTICAS Y FÍSICA (TUS PAQUETES ORIGINALES)
\usepackage{amsmath}
\usepackage{amssymb}
\usepackage{amsfonts}
\usepackage{amsthm}
\usepackage{bm}
\usepackage{physics}
\usepackage{fontawesome}
\usepackage{fp}
\usepackage{calculus}
\usepackage{calculator}
\usepackage{fancyvrb}
\usepackage{mathtools}
\usepackage{bigints}
\usepackage{fontawesome}

% PAQUETES DE GRÁFICOS E IMÁGENES (TUS PAQUETES ORIGINALES)
\usepackage{graphicx}
\usepackage{epstopdf}
\usepackage[export]{adjustbox}
\graphicspath{{./Fig}{./Img}{./Imgn}}

\usepackage[all]{nowidow}

% COLOR Y LAYOUT (MANTENIDOS)
\usepackage{xcolor}
\definecolor{orcidlogocol}{HTML}{A6CE39}
\usepackage{academicons}
\usepackage{tcolorbox}
\tcbuselibrary{breakable,skins}
\usepackage{multicol}
\usepackage{esvect}
\usepackage{relsize}
\usepackage{pifont}
\usepackage{etoolbox}
\newtcolorbox{graybox}{
	enhanced,
	boxrule=0pt,
	colback=gray!60!white,
	colupper=white,
	sharp corners,
	boxsep=5pt,
	left=1em, right=0.5em, top=0.5em, bottom=0.5em,
	nobeforeafter,
}

% FLOATS Y CAPTIONS (MANTENIDOS)
\usepackage{float}
\usepackage{placeins}
\usepackage{caption}
\usepackage{wrapfig}
\usepackage{array}
\usepackage{subcaption}
\usepackage{fancyvrb}
\usepackage[absolute,overlay]{textpos}

\captionsetup[table]{hypcap=false}
\captionsetup[figure]{hypcap=false}

%%%%%%%%%%%%%%%%%%%%%%%%%%%%%%%%%%%%%%%%%%%%%%%%%%%%%%%%%%%%%%%%%%%%%%%%%%%%%%%%%%%%%%%%%%%%%%%%%
%%% 9. CONFIGURACIONES Y COMANDOS (MANTENIDOS Y AJUSTADOS)
%%%%%%%%%%%%%%%%%%%%%%%%%%%%%%%%%%%%%%%%%%%%%%%%%%%%%%%%%%%%%%%%%%%%%%%%%%%%%%%%%%%%%%%%%%%%%%%%%

% CONFIGURACIÓN DE LEYENDAS (CAPTIONS)
\captionsetup[figure]{name=Figura, labelfont=bf, labelsep=period, font=footnotesize, justification=centering, width = 0.8\textwidth}
\captionsetup[table]{name=Tabla, labelfont=bf, labelsep=period, font=footnotesize, justification=centering, width = 0.8\textwidth}
\usepackage{enumitem}
\usepackage{listings}
\usepackage{needspace}




% COMANDOS PERSONALIZADOS Y ENTORNOS (TUS DEFINICIONES ORIGINALES)

\newcommand{\letraverbatim}{\fontsize{9pt}{11.5pt}\selectfont}
\newcommand{\ds}{\displaystyle}
\newcommand{\bi}{\int}
\newcommand{\la}{\langle}
\newcommand{\ra}{\rangle}
\newcommand{\ei}[1]{\mathbf{e}_{#1}}
\newcommand{\ptc}{\tcbset{colback=gray!10!white,colframe=black!85!blue}\begin{tcolorbox}[breakable,enhanced,title=\bf Parametrizando el sólido]}
\newcommand{\ptr}{\tcbset{colback=gray!10!white,colframe=black!85!blue}\begin{tcolorbox}[breakable,enhanced,title=\bf Parametrizando la región de Integración]}
\newcommand{\tc}{\begin{tcolorbox}[breakable,enhanced,title=\bf Instrucción en {\tt Mathematica}]}
\newcommand{\etc}{\end{tcolorbox}}
\newcommand{\ve}[1]{\mathbf{#1}}
\DeclareMathOperator{\arcsinh}{arcsinh}
\DeclareMathOperator{\arccosh}{arccosh}
\DeclareMathOperator{\arctanh}{arctanh}
\DeclareMathOperator{\arccoth}{arccoth}
\DeclareMathOperator{\arcsech}{arcsech}
\DeclareMathOperator{\arccsch}{arccsch}
\DeclareMathOperator{\rot}{rot}
\newcommand{\fig}{\begin{figure}}
\newcommand{\ef}[1]{\caption{#1}\end{figure}}
\newcommand{\vi}{\mathbf{i}}
\newcommand{\vj}{\mathbf{j}}
\newcommand{\vk}{\mathbf{k}}
\newcommand{\brake}{\\[4pt]} % Se ha corregido la sintaxis de este comando también% Comando auxiliar para saltos
		
% ENTORNOS DE TEOREMA Y REDEFINICIONES
\newtheorem{teorema}{Teorema}[section]
\newtheorem{prop}{Proposición}[section]
\newtheorem{defi}{Definición}[section]
\newtheorem{ejer}{Ejercicio}[section]
\newtheorem{ejemplo}{Ejemplo}[section]
\renewcommand{\proofname}{Demostración}
%\renewcommand{\proof}{Demostración} % Redefinición para Español
%Nueva redefinición al Español
\newcommand{\iniciodemostracion}{\noindent\textbf{Demostración.}\hspace{0.5em}}
\newcommand{\findemostracion}{\hfill$\square$}

\renewenvironment{proof}
{\vspace{0.5em}\iniciodemostracion\normalfont\ignorespaces} 
{\findemostracion\vspace{0.5em}\par} % 

\renewcommand*{\tablename}{\bf Tabla} % Estilo UD
\renewcommand*{\figurename}{\bf Figura} % Estilo UD
		
% AJUSTES DE ÍNDICE Y PROFUNDIDAD DE SECCIONES
\counterwithout{footnote}{chapter}
\setcounter{tocdepth}{3}
\setcounter{secnumdepth}{2} %-----> PONER NUMERO 3 O MÁS PARA CONTAR O NUMERAR
\renewcommand*\contentsname{\lhseries Contenido} % Estilo UD para Contenido
\renewcommand{\cfttoctitlefont}{\hfill\Huge}
\renewcommand{\cftdot}{}

% SOLUCIÓN DE CONFLICTOS
\AtBeginDocument{\RenewCommandCopy\qty\SI} % Solución de conflicto physics
\makeindex

\usepackage{fvextra}   % extiende fancyvrb y añade breaklines, breakanywhere
\fvset{breaklines=true, breakanywhere=true}

\titlespacing{\section}{0.75cm}{5mm}{2mm}
\titlespacing{\subsection}{0cm}{5mm}{2mm} 
\titlespacing{\subsubsection}{0pt}{5mm}{2mm}

\titleformat{\section}[block]{\color{MSBlack}\fontsize{16pt}{19pt} \fontssec}
{\thesection.\quad}{0pt}{\raggedright}

\titleformat{\subsection}{\color{MSBlack}\fontsize{13pt}{16pt}\fontssubsec}
{\thesubsection.\quad}{12pt}{}

%\titleformat{\subsubsection}[block]{\normalfont\fontsize{13pt}{2.0mm}}{}{0pt}{\thesubsubsection\quad} DESCOMENTAR EN CASO DE NUMERAR

\titleformat{\subsubsection}[block]
{\bfseries\fontsize{12pt}{14.5pt}\selectfont} % Negrita (estándar UD para sub-títulos)
{} % Sin etiqueta numérica
{0pt}
{\raggedright}

\renewcommand{\thesection}{\thechapter.\arabic{section}}
\renewcommand{\thesubsection}{\thesection.\arabic{subsection}}
\renewcommand{\thesubsubsection}{\thesubsection.\arabic{subsubsection}}

\raggedbottom
